\begin{papercard}{nouwens2020}
\cf{Stages} \stagetag{C} \stagetag{H} (also \stagetag{D} \stagetag{W}). Two studies of consent-management platforms (CMPs) after the GDPR. Study 1 scraped 680 CMP implementations on the UK top 10,000 sites and tested them against three \emph{legally derived, measurable} compliance conditions. Study 2 was a browser-extension field experiment ($n=40$) testing how eight designs change consent.

\cf{RQs} (1) What is the state of CMP interface design and how prevalent are non-compliant elements? (2) How do designs affect consent actions, i.e.\ how ``freely given'' consent is \pc{nouwens2020}{1}{What is the current state of interface design of CMPs in the EU}.

\cf{Legal criteria} Derived from the GDPR, the ePrivacy Directive, CJEU case law and regulator guidance. Implicit consent through continued use is invalid \pc{nouwens2020}{2}{So-called ‘implicit’ or ‘opt-out’ consent}. Pre-ticked boxes are invalid (recital 32). Emphasising ``agree'' over ``reject'' is non-compliant (ICO), and withdrawal must be as easy as giving consent (Art.\ 7(3)) \pc{nouwens2020}{2}{it must be “as easy to withdraw as to give consent”}. Purposes may not be bundled, and all vendors must be named \pc{nouwens2020}{2}{consent is invalid unless all organisations processing this data are specifically named}.

\cf{Study 1 method} A Scrapy and Splash scraper for the five largest third-party CMPs (about 58\,\% market share). It recorded the vendor, banner or barrier style, explicit vs.\ implicit consent and which actions count as consent, accept-all and reject-all buttons with the minimum clicks to reach them, and the lists, extent and defaults of purposes and vendors \pc{nouwens2020}{4}{The variables the scraper collected included the CMP vendor}. Run from a Danish IP over three days in September 2019 against the Alexa UK top 10k, throttled, rotating user agents \pc{nouwens2020}{4}{We ran the scraper from a Danish IP address}. The scraper was validated manually on 40 random sites per CMP \pc{nouwens2020}{5}{validated with 40 randomly selected sites from the list of 10,000}.

\cf{Compliance operationalisation} Three \emph{minimum} conditions testable with the scraped variables: consent is explicit; reject all is as easy as accept all (same number of clicks); no pre-ticked optional boxes. Meeting them does not guarantee compliance, so 11.8\,\% is a \emph{maximum} \pc{nouwens2020}{5}{This serves as a minimum hurdle: meeting these conditions alone will not guarantee compliance}. Excluded: qualitative information quality and visual prominence beyond clicks \pc{nouwens2020}{5}{Factors which could contribute to non-compliance which we did not examine}.

\cf{Study 1 results} 32.5\,\% used implicit consent (visiting, navigating, scrolling). 50.1\,\% had no reject-all button, and only 12.6\,\% offered reject-all with the same clicks as accept-all \pc{nouwens2020}{5}{Only 12.6\% of sites had a ‘reject all’ button accessible with the same or fewer number of clicks}. 56.2\,\% pre-ticked options. The median number of vendors was 315, and their descriptions averaged 7,985 words or about 31.9 minutes of reading \pc{nouwens2020}{6}{The mean total length of these descriptions per site is 7985 words}. Only \textbf{11.8\,\%} met all three conditions \pc{nouwens2020}{6}{Only 11.8\% of sites met these basic requirements}. Limitations: possible false positives and negatives in a dynamic DOM, customised and legacy CMPs, and anti-automation obfuscation \pc{nouwens2020}{6}{Although we manually validated the scraper, we cannot guarantee that there are no false negatives}.

\cf{Study 2 method} Two counterbalanced Latin-square within-subjects experiments. Experiment 1 ($2\times2$): notification style (barrier or banner) $\times$ bulk buttons (accept and reject, or accept only). Experiment 2 ($1\times4$): first-page granularity (bulk; +purposes; +vendors; +both). DV: the consent answer \pc{nouwens2020}{6}{Experiment 1 used a [2x2] latin square, within-subjects, repeated measures design}. A Chrome extension injected pop-ups into participants' \emph{normal browsing}, one every fourth URL, four repetitions per condition, with all interactions logged \pc{nouwens2020}{7}{The extension injected a pop-up every fourth url visited}. US residents were used to avoid real EU pop-ups, recruited by network and mailing list, paid \$50, and given a cover story about US vs.\ EU tracking \pc{nouwens2020}{7}{We selected this sample to prevent the confounding effects of real CMPs}. Eight of 48 were excluded for study-influenced answers. Analysis: linear regression with participant fixed effects \pc{nouwens2020}{8}{we created a linear regression model with fixed effects}.

\cf{Study 2 results} 89.3\,\% of answers used bulk buttons, and only 1.3\,\% were genuinely specific choices \pc{nouwens2020}{8}{only 17 answers out of 1280 (1.3\%) represent a participant consenting}. 93.1\,\% of interactions stayed on the first page: ``anything not immediately visible \ldots\ might as well not exist'' \pc{nouwens2020}{8}{anything not immediately visible to the user, anything requiring interaction to access, might as well not exist}. Style had no effect on acceptance, but banners were ignored 3.6 times more often. Removing reject-all from the first page raised consent by 22--23 points, and granular first pages lowered it by 8--20 points \pc{nouwens2020}{9}{meaning the accept rate increased on average by 23.1 percentage points}. Most participants said their answers did not match their ideal preferences, blaming interface design and fatigue \pc{nouwens2020}{9}{The most common reason mentioned, however, was the interface design}.

\cf{Limitations (Study 2)} A young, educated CS-adjacent sample (a ``best case''), uncontrolled exposure order, unequal exposures, and possible real pop-ups \pc{nouwens2020}{10}{There are a number of confounding variables that could have affected}.

\cf{Use for this thesis} Two transferable designs. (1) \emph{Operationalise legal requirements as observable, countable conditions}: explicit consent, reject as easy as accept (click counts), no pre-ticked options. Frame them as minimum hurdles giving an upper bound on compliance, which is exactly how a walkthrough codebook can check in-app consent without claiming legal conclusions. (2) Record \emph{click depth} to reject vs.\ accept, and whether choices are on the first layer. The finding that users stay on the first page justifies coding first-layer content separately. The experiment is out of scope, so cite it only for the effect sizes.
\end{papercard}
